\section{Abel summation and the pole at one}
\label{sec:abel-growth}

Write \(s=\sigma+it\), and let \(F\) be the entire extension of
\((s-1)\zeta(s)\), with \(F(1)=1\). The same function will be used at the
three heights \(0,t,2t\). For \(\sigma>0\), put
\[
  E(s)=\int_1^\infty \{u\}\,u^{-s-1}\,du,
  \qquad \{u\}=u-\lfloor u\rfloor .
\]
Here \(u^s=\exp(s\log u)\) for \(u>0\). We claim that
\begin{equation}
  F(s)=s\bigl(1-(s-1)E(s)\bigr)
  \qquad (\sigma>0).
  \label{eq:abel-entire}
\end{equation}

For \(\sigma>1\), the partial sums of the constant Dirichlet coefficient
\(1\) are \(\lfloor u\rfloor\). Abel summation therefore gives
\[
  \zeta(s)
  =s\int_1^\infty \lfloor u\rfloor u^{-s-1}\,du
  =s\left(\frac{1}{s-1}-E(s)\right).
\]
In the second equality we used \(\lfloor u\rfloor=u-\{u\}\) and
\(\int_1^\infty u^{-s}\,du=(s-1)^{-1}\). Both integrals in this
decomposition converge absolutely. In Lean the first equality is the
constant-coefficient specialization of \lean{LSeries_eq_mul_integral'}.

For the continuation, define \(h(u)=\mathbf{1}_{(1,\infty)}(u)\{u\}\)
on the positive real axis. Then
\(E(s)=\mathcal{M}h(-s)\), where
\(\mathcal{M}h(w)=\int_0^\infty u^{w-1}h(u)\,du\).
The measurable function \(h\) is bounded by one and vanishes in a
neighbourhood of zero. These estimates give holomorphy of the Mellin
transform for \(\Re w<0\) by
\lean{mellin_differentiableAt_of_isBigO_rpow}. Both sides of
\eqref{eq:abel-entire} are thus holomorphic on the connected half-plane
\(\sigma>0\), and the Abel identity gives equality on a neighbourhood of
\(s=2\). The identity theorem proves \eqref{eq:abel-entire} throughout the
half-plane, including at \(s=1\). These steps are formalized in
\lean{tail_smooth}, \lean{remainder_eq}, and \lean{entireFormula}.

Two consequences of the same remainder estimate will be needed. Absolute
integration gives
\begin{equation}
  |E(s)|\leq\int_1^\infty u^{-\sigma-1}\,du=\frac{1}{\sigma}.
  \label{eq:abel-remainder-bound}
\end{equation}
For \(\sigma\geq \tfrac12\), equations \eqref{eq:abel-entire}
and \eqref{eq:abel-remainder-bound} imply
\begin{equation}
  |F(s)|\leq |s|\bigl(1+2(|s|+1)\bigr)
           \leq (2|s|+2)^2 .
  \label{eq:quadratic-growth}
\end{equation}
These are \lean{mellin_bound} and \lean{growth} in the formal proof.

Suppose now that \(F(s)=0\) and \(\sigma\geq 3/4\).
Since \(F(1)=1\) and \(s\ne 0\), equation \eqref{eq:abel-entire} gives
\(E(s)=(s-1)^{-1}\). Hence \eqref{eq:abel-remainder-bound} and the triangle
inequality imply
\[
  \sigma \leq |s-1|\leq |\sigma-1|+|t|.
\]
This gives \(|t|\geq 2\sigma-1\) for \(3/4\leq\sigma\leq 1\)
and \(|t|\geq1\) for \(\sigma\geq1\). In particular,
\begin{equation}
  F(\sigma+it)=0,\quad \sigma\geq\tfrac34
  \quad\Longrightarrow\quad |t|\geq\tfrac12 .
  \label{eq:height-separation}
\end{equation}
This separation supplies the lower bound on the height used to estimate
the pole terms in the final argument.

\section{A local Landau inequality}
\label{sec:landau-disk}

Let \(D_r=\{w\in\mathbb C:|w|<r\}\). Suppose that \(f\) is entire,
\(f(0)\ne0\), every zero of \(f\) has nonpositive real part, and
\[
  |f(w)|\leq B|f(0)|\qquad (|w|\leq 1/2).
\]
In particular \(B\geq1\). Allow \(x\in\mathbb R\) either to be zero or to
satisfy \(x\leq-\Re(\rho^{-1})\) for some zero
\(\rho\in\overline D_{1/6}\). Then
\begin{equation}
  x-\Re\frac{f'(0)}{f(0)}\leq 24(\log B+1).
  \label{eq:landau-disk}
\end{equation}
The parameter \(x\) allows the same lemma to retain a specified zero at
one height and to bound the logarithmic derivative alone at the other
heights. Its two permitted cases are expressed by \lean{witness} and
the disjunction in \lean{landau}.

\subsection{Removing a zero}
\label{sec:landau-peel}

Choose a natural number \(n>2B\) and argue by induction on \(n\), allowing
\(f,B,x\) to vary. The case \(n=0\) is excluded by \(B\geq1\).
At a successor \(n=m+1\), suppose first that there is a zero \(\rho\)
as in the condition on \(x\). Since \(\rho\ne0\), the factorization
\(f(w)=(w-\rho)g(w)\) has \(g(0)\ne0\). The formal proof defines \(g\)
by the divided difference \lean{dslope}, whose value at \(\rho\) is \(f'(\rho)\).
Holomorphy of this divided difference follows from
\lean{differentiableOn_dslope}.

On \(|w|=1/2\), we have
\[
  |w-\rho|\geq \tfrac12-|\rho|\geq 2|\rho|,\qquad
  |f(0)|=|\rho|\,|g(0)|.
\]
Consequently \(|g(w)|\leq (B/2)|g(0)|\) on the circle. The maximum
modulus principle extends the estimate to its closed disk. All zeros of
\(g\) are zeros of \(f\). The induction hypothesis applies to \(g\) with
bound \(B/2\) and with \(x=0\), since \(2B<m+1\) and \(B\geq1\) give
\(2(B/2)=B<m\). Logarithmic differentiation gives
\[
 \frac{f'(0)}{f(0)}=-\rho^{-1}+\frac{g'(0)}{g(0)},\qquad
 x-\Re\frac{f'(0)}{f(0)}
   \leq -\Re\frac{g'(0)}{g(0)}
   \leq 24\bigl(\log(B/2)+1\bigr).
\]
Monotonicity of the real logarithm now gives \eqref{eq:landau-disk}.
The Lean source credits this monotonicity step after division of the
growth bound to Pi's public Core~1 submission \#96263.

\subsection{The zero-free disk}
\label{sec:landau-regular}

In the remaining branch the hypothesis on \(x\) forces \(x=0\).
Any zero \(\rho\in D_{1/6}\) would satisfy
\(-\Re(\rho^{-1})=-\Re(\rho)/|\rho|^2\geq0\), supplying a witness even
for \(x=0\). Thus \(f'/f\) is holomorphic on \(D_{1/6}\). It has a
primitive \(G\) there with \(G(0)=0\). The logarithmic derivatives of
\(f\) and \(\exp G\) agree, so their quotient is constant on the disk.
Evaluation at the origin gives \(f(w)=f(0)\exp G(w)\). The formal proof
uses \lean{isExactOn_ball} for the primitive and \lean{logDeriv_eqOn_iff}
for this identity.

Set \(A=\log B+1>0\). The original bound gives
\(\Re G(w)=\log(|f(w)|/|f(0)|)\leq\log B\) on \(D_{1/6}\).
The Borel--Carath\'eodory inequality for a function vanishing at the
centre gives
\[
  |G(w)|\leq \frac{2A|w|}{1/6-|w|}\qquad (|w|<1/6).
\]
In particular, \(|G(w)|\leq 2A\) on \(|w|=1/12\). Cauchy's derivative
estimate on that circle yields
\[
  -\Re\frac{f'(0)}{f(0)}\leq |G'(0)|
                           \leq \frac{2A}{1/12}
                           =24(\log B+1).
\]
This proves \eqref{eq:landau-disk} in the second branch. The two estimates
are implemented by \lean{borelCaratheodory_zero} and
\lean{norm_deriv_le_of_forall_mem_sphere_norm_le}.
